\documentclass[10pt,a4paper]{amsart}
\usepackage[margin=19mm]{geometry}
\usepackage{amsmath,amssymb,mathtools}
\usepackage[scr=rsfs,cal=euler]{mathalfa}
\usepackage{xfrac}
\usepackage[unicode,hidelinks,bookmarks=false]{hyperref}
\newcommand{\ie}{i.e.\ }
\newcommand{\cf}{cf.\ }
\newcommand{\resp}{respectively\ }
\newcommand{\Subsection}{\subsection}
\providecommand{\nobreakdash}{\nobreak\hbox{-}}
\setlength{\emergencystretch}{2em}
\multlinegap0pt
\usepackage{mathtools}
\usepackage[all]{xy}
\usepackage{amssymb}
\usepackage{enumerate}
\usepackage{xcolor}
\usepackage{algorithm}\usepackage{algpseudocode}
\usepackage{xfrac}
\usepackage{tikz}
\newtheorem{lem}{Lemma}
\newtheorem{claim}{Claim}
\newcommand{\U}{X}
\newcommand{\diam }[1]{{\textbf{diam}\big(#1\big)}}
\title{Growth tightness and genericity for word metrics from injective spaces}
\author{Lihuang Ding, D\'idac Mart\'inez-Granado, Abdul Zalloum}
\date{Source: arXiv:2407.02378v4 (CC BY 4.0)}
\begin{document}
\maketitle

\noindent\emph{Source and licence.} Growth tightness and genericity for word metrics from injective spaces. Version arXiv:2407.02378v4 (CC BY 4.0), distributed by arXiv under the Creative Commons Attribution 4.0 International licence, http://creativecommons.org/licenses/by/4.0/, as recorded in the arXiv licence metadata for this exact version and shown on the abstract page https://arxiv.org/abs/2407.02378v4. Full licence notice: \texttt{licenses/arxiv-2407.02378-CC-BY-4.0.txt}.\par\smallskip

\noindent\textit{Editorial scope.} Counted unit: the article's lem projLip (source0.tex lines 1208--1211). The article's complete original proof of it is delivered with it (source0.tex lines 1212--1247, one \texttt{proof} environment of 36 lines). No mathematics is written, repaired or re-derived: the statement and the proof are the article's own lines, in the article's own notation, and no retained line is substituted or re-flowed.  Retained uncounted as the source-local context the proof needs: lines 1102--1116.  Nothing was omitted from any retained span, so the package carries no omission record.  No retained line contains a citation, so the wrapper carries no bibliography and no bibliography entry.  No retained line contains a figure environment, an image-inclusion command or drawing code, so no image is delivered and the package declares no asset dependency.  Editorial formatting only, in addition: an AMS \texttt{amsart} preamble that restates the article's own declarations and macros used by the retained lines, namely the environments \texttt{lem}, \texttt{claim} on separate counters, together with the standard packages those lines need.  The retained environments print in this document as Lemma 1, Claim 1 in order of appearance, whereas the article prints the same items with the numbering of its own full document.  The complete native source of the article as distributed in the arXiv e-print for this exact version is bundled unchanged as \texttt{sources/161714/source0.tex}.
\begin{lem}\label{BigFive}
	Let $Y\subseteq \U$ be a  $C$-contracting subset for $C>0$. Then the following holds.
	
	\begin{enumerate}
		\item
		For any  geodesic $\gamma$, we have $$\big |d^{\pi}_Y(\gamma^-,\gamma^+)- \diam{\gamma\cap N_C(Y)}\big|\leq 4C.$$
		
		\item
		There exists $\hat C=\hat C(C)$ such that   $d^{\pi}_Y(y, z)\le d(y, z)+\hat C$ for any $y,z\in \U$.
		
		\item 
		There exists $\hat C=\hat C(C)$ such that for any geodesic $\gamma$, we have $$\left| \diam{\pi_Y(\gamma)} - d^{\pi}_Y(\gamma_-, \gamma_+) \right| \leq \hat C$$
		
	\end{enumerate}
\end{lem}

\begin{lem}\label{projLip}
	For any $\beta>0$, $C\ge 0, K\ge 1$, there exist $a = a(K, \beta) \ge 0$, $b = b(\beta, K, C) > 0$ and $c = c(K) > 0$ with the following property. 
	Let $Y$ be a $\beta$-contracting subset of $X$ and $\gamma$ be a $(K, C)$-quasi-geodesic with endpoints $x\in X$ and $y\in Y$. Then $$\diam {\pi_{Y}(\gamma)} \le a \log d(x, y) + b + c d(\pi_Y(x), y)).$$ Furthermore,  we may choose $a(1, \beta)=0$ and $\lim_{K\to 1} a(K, \beta) = 0$ for any $\beta > 0$. 
\end{lem}

\begin{proof}
	Set $a(1, \beta)=0$ and assume that $K>1$. Denote $\theta \coloneqq 1-\frac{1}{K}$.
	Let $z$ be the first point on $\gamma$ such that $d(z, Y) \le  \theta d(x, Y)$ and let $\gamma_1$ be the subpath of $\gamma$ between $x$ and $z$. Let $\mathfrak h = \max\left\{\frac{2K^2\beta}{K-1}, \frac{CK}{K^2-1}, 1\right\}$.
	\begin{claim}
		If $d(x, Y) > \mathfrak h$, then $\diam{\pi_Y(\gamma_1)} \le 6(K + 1) \beta$.
	\end{claim}
	\begin{proof}[Proof of the claim]
		Let $D = \diam{\pi_Y(\gamma_1)}$, $l = \mathrm{length}(\gamma_1)$ and $h = d(x, Y)$. Then by triangular inequality, $$d(x, z) \le d(x, Y) + d(z, Y) + d^{\pi}_Y(x, z) \le h + \theta h + D.$$ 
		Since $\gamma$ is a $(K, C)$-quasi-geodesic, we have \begin{align}\label{lboundEq}
			l \le K d(x, z) + C \le K(h(1+\theta) + D) + C
		\end{align}
		
		We are going to bound $D$. To that end,  we write $\gamma_1 = \gamma_2 \cup \gamma_3$ as the concatenation of two paths where  the subpath $\gamma_2$ has length $h$ starting at $x$ and thus is contained in $B(x, h)$, and $\gamma_3$ is the remaining subpath of $\gamma_1$ with $\mathrm{length}(\gamma_3) = l - h$. Since $Y$ is $\beta$-contracting and $h = d(x, Y)$, we have $\diam{\pi_Y(\gamma_2)} \le \beta$. For $\gamma_3$, we may divide it into disjoint subpaths of length $\theta h$, and the number of them is at most $({l-h})/{\theta h} + 1.$ Since the path $\gamma_1$ from $x$ to $z$ has distance at least $\theta h$ to $Y$, we have that each of these subpaths is contained in a ball which is disjoint from $Y$. Thus, $$\diam{\pi_Y(\gamma_3)}\le \beta\left({(l-h)}/{\theta h} + 1\right).$$ Taking into account the projections  of $\gamma_2$ and $\gamma_3$, we have $$D = \diam{\pi_Y(\gamma_1)} \le \diam{\pi_Y(\gamma_2)} + \diam{\pi_Y(\gamma_3)} \le \beta \left(2 + {(l-h)}/{\theta h}\right)$$
		which yields $$l \ge h + \left({D}/{\beta} - 2\right)\theta h.$$
		With (\ref{lboundEq}) we have $$h + \left({D}/{\beta} - 2\right) \theta h \le K\left(h(1+\theta) + D\right) + C,$$ or equivalently, using  $\theta=1-\frac{1}{K}$, $$D\left( {\theta h}/{\beta} - K\right)  \le h\left(K +K\theta- 1+2\theta \right) + C =  2h(K +1 )\theta + C.$$ Recall that   $\mathfrak h= \max\{\frac{2K\beta}{\theta}, \frac{C}{(K+1)\theta}, 1\}$. If $h > \mathfrak h$, we have ${h\theta }/{\beta} \ge 2K$ and $C\le h (K +1) \theta$. Thus, $$D \le \frac{3h(K +1)\theta}{{\theta h}/{2\beta}} = 6(K + 1)\beta.$$
		This proves the claim.
	\end{proof}
	Let $x_0 = x$ and define inductively $x_k$ to be the first point on $\gamma$ with $d(x_k, Y) \le \theta d(x_{k-1}, Y)$ for $k\ge 1$. Let $m$ be the minimal $k$ such that $d(x_k, Y)\le \mathfrak h$. Then $$-{\log_{\theta} \frac{d(x,Y)}{\mathfrak h}} \le m \le 1- {\log_\theta \frac{d(x,Y)}{\mathfrak h}} \le 1 - \log_\theta d(x,Y).$$  
	Let $\gamma_k$ be the subpath of $\gamma$ between $x_{k-1}$ and $x_k$. As    $d(\gamma_{k}, Y)\ge \mathfrak h$ for $k\le m$, by the claim above, we have $\diam{\pi_Y(\gamma_k)}\le 6(K+1) \beta$. If $\alpha_1$ denotes the subpath of $\gamma$ between $x$ and $x_m$, then $$\diam{\pi_Y(\alpha_1)} \le 6m(K+1) C,$$ and especially $d_Y^\pi(x, x_m) \le 6m(K+1) \beta$. Thus, $d(x_m, y)\le d(x_m, \pi_Y(x_m)) + d_Y^\pi(x, x_m) + d(\pi_Y(x), y) \le 6m(K+1)\beta + \mathfrak h + d(\pi_Y(x), y)$. Let $\alpha_2$ be the subpath of $\gamma$ between $x_m$ and $y$. Since $\alpha_2$ is also a $(K, C)$-quasi-geodesic, we have $\mathrm{length}(\alpha_2)\le K d(x_m, y) + C$. So by Lemma \ref{BigFive}, $$\diam{\pi_Y(\alpha_2)} \le \mathrm{length}(\alpha_2) + \hat C \le 6m K(K+1) \beta + K \mathfrak h + K d(\pi_Y(x), y) + C + \hat C.$$ Finally, $$\begin{aligned}
		\diam{\pi_Y(\gamma)} &\le \diam{\pi_Y(\alpha_1)} + \diam{\pi_Y(\alpha_2)}\\
		&\le 6m (K+1)^2 \beta + K \mathfrak h + K d(\pi_Y(x), y) + C + \hat C\\
		&\le a \log d(x, y)  +b + c d(\pi_Y(x), y) 
	\end{aligned}$$
	where $a = -{6(K+1)^2 \beta}/{\log \theta}$, $b = {6(K+1)^2 \beta} + K \mathfrak h + C + \hat C$ and $c = K$. 

    \iffalse
    Furthermore, let $z\in \gamma$. If $z\in \alpha$, then $d(z, y) \le \mathrm{length}(\alpha) \le 6mK(K+1) C + C$ and $\log(d(z, y)) < d(z, y) \le K(a_0 \log d(x, y)  +b_0 + c_0 d(\pi_Y(x), y)) + C$. Thus, $$\begin{aligned}
        d^\pi(z, y)&\le \diam {\pi_Y(\gamma)} \le a_0 \log d(x, y) + b_0 + c_0 d(\pi_Y(x), y)\\
        &\le a_0 \log \frac{d(x, y)}{d(z, y)} + b_0 + c_0 d(\pi_Y(x), y) + a_0 \log {d(z, y)}\\
        &\le (1+K a_0) (a_0 \log d(x, y)  +b_0 + c_0 d(\pi_Y(x), y) + C
    \end{aligned}$$
    If $z\in \beta$, then there exists $1\le k\le m$ such that $z\in \gamma_k$. Thus, $d^{\pi}_Y(z, x) \le 6(K+1) C k$. We have $d(z, y) \le \mathrm{length}([x_{k-1}, y]_\gamma) \le K d(x_{k-1}, y) + C \le K(d^{\pi} (x_{k-1}, y) + d(x_{k-1}, Y))$. Thus $k\le -\log_{\theta} \frac{d(x, Y)}{d(z, y)}$
    \fi
    
    Noticing  $\lim_{K\to 1^+} a = 0$, the lemma thus follows.
\end{proof}

\end{document}
